\chapter*{ЗАКЛЮЧЕНИЕ}
\addcontentsline{toc}{chapter}{ЗАКЛЮЧЕНИЕ}

В данной работе разработана и экспериментально проверена самообучающаяся модель искусственного интеллекта для синтеза и оценки мелодий по биологическим данным.

\section*{Основные результаты}

\textbf{1. Анализ предметной области.} Изучены существующие подходы к биосонификации и генерации символической музыки. Выявлены недостатки: прямое отображение символов даёт музыкально бедный результат, алгоритмические методы требуют ручной настройки, генеративные модели без структурного разделения создают хаотичные последовательности.

\textbf{2. Иерархическая архитектура.} Разработана архитектура Bio→Harmony→Melody, где биологические признаки используются как управляющий сигнал для структурированной музыкальной генерации. Система разделяет задачу на два этапа: генерацию гармонии и генерацию мелодии поверх гармонии.

\textbf{3. Биологический энкодер.} Реализовано извлечение признаков из DNA, RNA и protein последовательностей: состав оснований, энтропия, GC-content, белковые свойства (ароматичность, индекс нестабильности, изоэлектрическая точка, GRAVY). Энкодер возвращает 256-мерный embedding и 6-мерный control profile.

\textbf{4. Музыкальное представление.} Разработано структурированное представление музыки с разделением на гармонию (HarmonyBar) и мелодию (MelodyEvent). Реализованы токенизаторы для кодирования аккордовых последовательностей и мелодических событий.

\textbf{5. Pairing модуль.} Реализовано сопоставление биологических фрагментов с музыкальными сегментами через пространство дескрипторов (темп, плотность, частота смены аккордов, регистр, сложность, модальность).

\textbf{6. Нейросетевая модель.} Реализован Bio-Conditioned Transformer. Актуальная модель веб-интерфейса использует параметры: d\_model=384, n\_layers=6, n\_heads=6, dim\_feedforward=1536, dropout=0.20. Модель использует bio conditioning через memory tokens.

\textbf{7. Обучение.} Система обучена на корпусе POP909 и геномных данных RefSeq. Актуальная модель использовала 3746 биологических фрагментов, 25940 музыкальных сегментов и 15925 обучающих пар. Обучение выполнено на GPU RTX 2060 6GB с использованием mixed precision. Модель показала validation loss 0.1454 для гармонии и 0.1566 для мелодии.

\textbf{8. Экспериментальная оценка.} Проведено сравнение актуальной web-модели с предыдущей дипломной моделью и random baseline. Актуальная модель показала chord-tone ratio 0.6710 против 0.6073 у предыдущей модели и 0.4043 у baseline, что подтверждает улучшение согласованности мелодии и гармонии.

\textbf{9. Воспроизводимость.} Все результаты воспроизводимы: фиксированный seed, сохранённые checkpoint, документированные конфиги, отчёты по данным и метрикам.

\section*{Достижение цели и решение задач}

Цель работы достигнута: разработана и экспериментально проверена самообучающаяся модель искусственного интеллекта для синтеза и оценки мелодий по биологическим данным.

Все поставленные задачи решены:
\begin{itemize}
\item Изучена предметная область биосонификации и генерации символической музыки
\item Проанализированы существующие решения и выявлены их недостатки
\item Подготовлены биологические данные (FASTA) и музыкальный корпус (MIDI)
\item Реализовано извлечение биологических признаков из DNA, RNA и protein последовательностей
\item Разработана иерархическая архитектура генерации: Bio→Harmony→Melody
\item Реализован pairing биологических и музыкальных сегментов
\item Обучены модели на доступном оборудовании (NVIDIA RTX 2060 6GB)
\item Проведена экспериментальная оценка и сравнение с baseline
\item Подготовлены воспроизводимые результаты и документация
\end{itemize}

\section*{Ограничения}

Разработанное решение имеет ограничения:
\begin{itemize}
\item Модель зависит от стиля корпуса POP909
\item Ограничение GPU памяти не позволило использовать ESM embeddings
\item Веб-генерация строится из 4-тактовых фрагментов и не формирует полноценную длинную музыкальную форму
\item Система не доказывает причинную связь между биологией и музыкой
\item Музыка зависит от стиля обучающего корпуса (POP909)
\item Мелодия строго монофоническая
\item Гармония ограничена 11 качествами аккордов
\end{itemize}

\section*{Практическая значимость}

Создан рабочий программный прототип BioSonification, включающий:
\begin{itemize}
\item Обученные модели с воспроизводимыми checkpoint
\item Веб-интерфейс для генерации музыки из FASTA
\item Документацию и примеры использования
\item Отчёты по данным и метрикам качества
\end{itemize}

Система может использоваться в образовательных целях, для научной коммуникации и демонстрации биоинформатических данных.

\section*{Направления дальнейшего развития}

Возможные направления развития системы:
\begin{enumerate}
\item \textbf{Увеличение корпуса.} Обучение на большем музыкальном корпусе (MAESTRO, Lakh MIDI) для повышения разнообразия.

\item \textbf{Использование ESM embeddings.} Интеграция предобученных белковых embeddings для более богатого биологического представления.

\item \textbf{Расширение длины генерации.} Генерация композиций длиннее 8 тактов с сохранением музыкальной связности.

\item \textbf{Добавление аккомпанемента.} Генерация третьей дорожки с ритмическим или басовым аккомпанементом.

\item \textbf{Динамика и артикуляция.} Добавление velocity и артикуляционных событий для более выразительной музыки.

\item \textbf{Интерактивный интерфейс.} Разработка интерактивного веб-приложения с визуализацией биологических признаков и музыкальных параметров.

\item \textbf{Оценка музыкального качества.} Проведение перцептивного эксперимента с участием слушателей для оценки музыкальной выразительности.

\item \textbf{Применение к другим типам данных.} Адаптация системы для сонификации других научных данных (временные ряды, графы, изображения).
\end{enumerate}

\section*{Заключительные замечания}

Разработанная система демонстрирует возможность использования биологических признаков как управляющего сигнала для генерации музыкально связных композиций. Иерархическая архитектура обеспечивает структурную организацию музыки, а обучение на реальном музыкальном корпусе — музыкальную выразительность.

Результаты работы подтверждают, что сонификация может быть не только прямым отображением символов, но и управляемой генерацией, создающей музыку, которая отражает свойства биологических данных через музыкальные параметры. Повторная оценка показала, что обновлённая web-модель улучшает техническую согласованность гармонии и мелодии по сравнению как со случайным baseline, так и с предыдущей версией модели.
